% E0-A CAL-only FP64 pilot: fixed protocol, 80 solver evaluations.
% Figures here transcribe the completed host audit output. Raw archive
% is pending the separate research-branch upload as of this commit.
\section{FP64 solver operating-point calibration (E0)}
\label{sec:e0-calibration}

The original 50,000-update linear inverse-temperature schedule
was used in M5 and M6 as a \emph{finite-budget operating point},
not as a claim of convergence to the published Max-Cut optimum.
A preregistered CAL-only pilot therefore compared the same
legacy linear schedule with a geometric annealing schedule
calibrated from conditional \emph{uphill} acceptance probabilities
on reference-weight random-spin states. Its temperature endpoints
were fixed before any E0 solver execution; held-out instances
were not accessed.

The CAL graphs were \texttt{ising2.5-100\_5555} and
\texttt{t2g10\_5555}, each with 100 spins. At each graph,
20 exactly matched FP64 seeds were run for 50,000 attempted
spin updates with either the historical linear schedule
or the newly frozen geometric schedule. Thus E0 comprised
$2\times2\times20=80$ solver evaluations. The on-host
read-only audit passed for all 80 unique run records,
320 checkpoint entries, 800 acceptance-trace bins,
16 quality summary rows, eight paired difference
rows, source and protocol hashes, and original optimum
evaluation. It launched no further solver runs.

\begin{table}[htbp]\centering\small
\begin{tabular}{@{}l rrr@{}}\toprule
CAL instance & Legacy mean gap (\%) &
Geometric mean gap (\%) & Geometric--linear (pp)\\\midrule
Dense & 4.713 & 4.388 & $-0.325$\\
Sparse & 3.607 & 3.887 & $+0.280$\\\bottomrule
\end{tabular}
\caption{E0 original-weight optimum-gap means at 50,000
attempted updates per run; 20 paired seeds per schedule and
graph. Lower gap is better. Both schedule differences require
their paired confidence intervals for inference.}
\label{tab:e0-means}
\end{table}

The final 50,000-update paired geometric-minus-linear
95\% seed-bootstrap confidence interval is
$[-1.183,+0.533]$ pp on the dense CAL and
$[-0.566,+1.084]$ pp on the sparse CAL. Both
include zero. The equal-graph average mean gap is
4.160\% for legacy linear and 4.138\% for geometric.
These descriptive averages do not establish an
overall schedule advantage. Neither schedule on
either graph attained the exploratory 1\% mean
optimum-gap target.

The trace audit explains why schedule shape mattered.
The geometric schedule was initially very hot,
with first-5000-update observed spin-flip rates
43.66\% (dense) and 44.00\% (sparse), but
its last-bin flip rates fell to only 0.10\%
and 0.15\%, respectively. The legacy linear
last-bin rates were 2.34\% and 8.33\%.
Similarly the geometric last-bin conditional
uphill acceptance rates were 0.05\% (dense)
and 0.08\% (sparse), compared with 1.19\%
and 4.41\% under linear. These time-resolved
data demonstrate strong late cooling but are
not proof of convergence or a unique causal
explanation for outcome quality.

On sparse CAL, geometric had a matched-seed
gap advantage of $-3.127$ pp at 30,000 updates
(95\% CI $[-4.342,-1.922]$), but at 50,000
updates the paired contrast was $+0.280$ pp
($[-0.566,+1.084]$), reflecting the fixed-budget
comparison's time dependence. Both graph
final-contrast intervals include zero.

\noindent\textbf{E0 decision.}
The pilot does not support adopting the geometric
schedule as the new reference. The original
50,000-update \emph{linear} schedule will be
retained as a reproducible, not-proven-converged
finite-budget comparison condition for the
\emph{prospective} minimum-width study, subject
to supervisor confirmation. The E0 1\% optimum-gap
aspiration is separate from the proposed $+0.5$
percentage-point \emph{additional gap} tolerance
used to certify new low-bit formats against
same-budget, same-response FP64. No existing M5/M6
results or quality screens have been revised.

\noindent\textbf{Provenance.}
The source-frozen E0 protocol and read-only audit
utility are in the
\href{https://github.com/Makise-okabe/MRAM_P-bit/tree/research/m4-gain-fusion-20261107}
{research branch}. This chapter is grounded in
the successful on-host integrity-audit console
record; the full E0 raw CSV/manifest archive is now
retained in the
\href{https://github.com/Makise-okabe/MRAM_P-bit/tree/research/m4-gain-fusion-20261107/results/e0_fp64_two_cal_20261112}
{research-branch E0 results directory}, with 80 run
rows, 320 checkpoints, 800 trace bins and the
frozen protocol and manifest.
